% Problem record op_b315f0d0b6ddbdee. This TeX file is derived from
% database/problems_json/op_b315f0d0b6ddbdee.json by scripts/sync-tex.mjs; edit the JSON record
% and rerun the script rather than editing this file.
\section{Minimal dimensions for strict transpose degradability}
\label{sec:op_b315f0d0b6ddbdee}

\paragraph{Problem.}
What are the componentwise-minimal dimension triples
$(d_A,d_B,d_E)$ that admit a transpose-degradable but nondegradable channel?
Let $d_X:=\dim X$ and let $V:A\to B\otimes E$ be a support-minimal isometry,
meaning that the channels
\begin{equation}
  \Phi_V(X):=\operatorname{Tr}_E(VXV^\dagger),
  \qquad
  \Phi_V^c(X):=\operatorname{Tr}_B(VXV^\dagger)
  \label{eq:p54-complementary-pair}
\end{equation}
satisfy
$\operatorname{supp}(\Phi_V(I_A))=B$ and
$\operatorname{supp}(\Phi_V^c(I_A))=E$.
Equation~\eqref{eq:p54-complementary-pair} is strictly transpose degradable
when, for a fixed-basis transpose $\mathsf T_E$, its factorization properties
are
\begin{equation}
  \begin{aligned}
    &\exists\ \mathcal D:\mathcal L(B)\to\mathcal L(E)\ \text{CPTP},
    &&\mathsf T_E\circ\Phi_V^c=\mathcal D\circ\Phi_V,\\
    &\nexists\ \widetilde{\mathcal D}:\mathcal L(B)\to\mathcal L(E)\
      \text{CPTP},
    &&\Phi_V^c=\widetilde{\mathcal D}\circ\Phi_V.
  \end{aligned}
  \label{eq:p54-strict-factorizations}
\end{equation}
A feasible triple is componentwise minimal if no distinct feasible
$(d'_A,d'_B,d'_E)$ satisfies $d'_X\leq d_X$ for every $X\in\{A,B,E\}$.
Determine all minimal triples satisfying
Eq.~\eqref{eq:p54-strict-factorizations}.

\subsection*{Status}

Unsolved

\subsection*{Source}

This dimension-refined problem is implicit in Singh and Datta's explicit
question about whether transpose degradability differs from degradability and
in their support-minimal dimension formalism
\sourcecite{ref:p54-singh-datta}{SD22}.

\subsection*{Progress}

\begin{itemize}
  \item Support minimality in Eq.~\eqref{eq:p54-complementary-pair} gives the
  Choi-rank identities
  \begin{equation}
    d_E=\operatorname{rank}J(\Phi_V),
    \qquad
    d_B=\operatorname{rank}J(\Phi_V^c).
    \label{eq:p54-choi-ranks}
  \end{equation}
  Equation~\eqref{eq:p54-choi-ranks} removes artificial output or environment
  dimensions introduced by nonminimal dilations
  \sourcecite{ref:p54-singh-datta}{SD22}.

  \item The first line of Eq.~\eqref{eq:p54-strict-factorizations} makes
  $J(\Phi_V^c)$ PPT, whereas separability of this Choi operator would make
  $\Phi_V$ degradable.  The low-rank PPT separability theorem therefore gives
  the necessary inequality
  \begin{equation}
    d_B=\operatorname{rank}J(\Phi_V^c)
      >\max\{d_A,d_E\}.
    \label{eq:p54-rank-obstruction}
  \end{equation}
  Equation~\eqref{eq:p54-rank-obstruction} is only an obstruction: a
  PPT-entangled Choi operator need not satisfy the channel factorization in
  Eq.~\eqref{eq:p54-strict-factorizations}
  \sourcecite{ref:p54-bradler}{Bra15},
  \sourcecite{ref:p54-horodecki-low-rank}{HLVC00}.

  \item PPT is equivalent to separability on $2\otimes2$ and $2\otimes3$.
  Combining this fact with Eq.~\eqref{eq:p54-rank-obstruction}, the
  componentwise-minimal triples not excluded by the known tests are
  \begin{equation}
    (d_A,d_B,d_E)\in
      \{(2,5,4),\ (3,4,3),\ (4,5,2)\}.
    \label{eq:p54-candidate-triples}
  \end{equation}
  No triple in Eq.~\eqref{eq:p54-candidate-triples} is known to be realizable
  \sourcecite{ref:p54-horodecki-ppt}{HHH96},
  \sourcecite{ref:p54-horodecki-low-rank}{HLVC00}.
\end{itemize}

\subsection*{References}

\begin{enumerate}[label={},leftmargin=4.8em,labelsep=0.5em]
  \footnotesize
  \item[\textup{[SD22]}]\label{ref:p54-singh-datta}
  S. Singh and N. Datta,
  \lq\lq Detecting Positive Quantum Capacities of Quantum Channels,\rq\rq{}
  \emph{npj Quantum Information} \textbf{8}, 50 (2022).
  \href{https://doi.org/10.1038/s41534-022-00550-2}{doi:10.1038/s41534-022-00550-2};
  \href{https://arxiv.org/abs/2105.06327}{arXiv:2105.06327}.

  \item[\textup{[Bra15]}]\label{ref:p54-bradler}
  K. Br{\'a}dler,
  \lq\lq The Pitfalls of Deciding Whether a Quantum Channel Is (Conjugate)
  Degradable and How to Avoid Them,\rq\rq{}
  \emph{Open Systems \& Information Dynamics} \textbf{22}, 1550026 (2015).
  \href{https://doi.org/10.1142/S1230161215500262}{doi:10.1142/S1230161215500262};
  \href{https://arxiv.org/abs/1507.06159}{arXiv:1507.06159}.

  \item[\textup{[HLVC00]}]\label{ref:p54-horodecki-low-rank}
  P. Horodecki, M. Lewenstein, G. Vidal, and I. Cirac,
  \lq\lq Operational Criterion and Constructive Checks for the Separability of
  Low-Rank Density Matrices,\rq\rq{} \emph{Physical Review A} \textbf{62},
  032310 (2000).
  \href{https://doi.org/10.1103/PhysRevA.62.032310}{doi:10.1103/PhysRevA.62.032310};
  \href{https://arxiv.org/abs/quant-ph/0002089}{arXiv:quant-ph/0002089}.

  \item[\textup{[HHH96]}]\label{ref:p54-horodecki-ppt}
  M. Horodecki, P. Horodecki, and R. Horodecki,
  \lq\lq Separability of Mixed States: Necessary and Sufficient
  Conditions,\rq\rq{} \emph{Physics Letters A} \textbf{223}, 1--8 (1996).
  \href{https://doi.org/10.1016/S0375-9601(96)00706-2}{doi:10.1016/S0375-9601(96)00706-2};
  \href{https://arxiv.org/abs/quant-ph/9605038}{arXiv:quant-ph/9605038}.
\end{enumerate}

\subsection*{Comment}

The feasible set may be empty because existence of any strict
transpose-degradable channel is unresolved in Problem~53.  Under channel
complementation, the triples for the equivalent strict
transpose-antidegradable formulation are obtained by interchanging $d_B$ and
$d_E$.

\subsection*{Field}

Quantum Shannon theory

\subsection*{Topic}

Transpose-degradable channels; Degradable channels; Partial transpose criterion; Choi states; Quantum channels

\subsection*{ID}

\texttt{op\_b315f0d0b6ddbdee}
